\documentclass[10pt,a4paper]{amsart}
\usepackage[margin=19mm]{geometry}
\usepackage{amsmath,amssymb,mathtools}
\usepackage[scr=rsfs,cal=euler]{mathalfa}
\usepackage{xfrac}
\usepackage[unicode,hidelinks,bookmarks=false]{hyperref}
\newcommand{\ie}{i.e.\ }
\newcommand{\cf}{cf.\ }
\newcommand{\resp}{respectively\ }
\newcommand{\Subsection}{\subsection}
\providecommand{\nobreakdash}{\nobreak\hbox{-}}
\setlength{\emergencystretch}{2em}
\multlinegap0pt
\usepackage{bm}
\usepackage{bbm}
\usepackage{dsfont}
\usepackage{xspace}
\usepackage{cancel}
\usepackage{mathtools}
\usepackage{amsthm}
\makeatletter
\@ifundefined{thm}{\newtheorem{thm}{Theorem}}{}
\makeatother
\title[Classification and Ideal Lattices of Leavitt Path Algebras]{Classification and Ideal Lattices of Leavitt Path Algebras}
\author{Yvan Grinspan, Seth Yoo}
\date{Source: arXiv:2509.25464v1 (CC BY 4.0)}
\begin{document}
\maketitle

\noindent\emph{Source and licence.} Classification and Ideal Lattices of Leavitt Path Algebras. Version arXiv:2509.25464v1 (CC BY 4.0), distributed under the Creative Commons Attribution 4.0 International licence, as recorded in the arXiv licence metadata for this exact version and shown on the abstract page https://arxiv.org/abs/2509.25464v1. Full rights notice: licenses/arxiv-2509.25464-CC-BY-4.0.txt.\par\smallskip

\noindent\textit{Editorial scope.} Counted unit: the Thm whose source label is \texttt{None} in the article (source0.tex lines 396--398, source label \texttt{None}), delivered together with the article's own complete original proof of it (source0.tex lines 399--441, one \texttt{proof} environment of 43 lines). No mathematics is written, repaired or re-derived: statement and proof are the article's own text line for line. Required context retained uncounted: none. Every definition, display and label of those lines is retained unchanged. Omitted from this package and not redistributed: 1--395, 442--572. Nothing inside a retained excerpt is omitted or altered: every retained body is delivered exactly as the article's lines are written, so no omission receipt of any kind accompanies this package. Citations: the retained excerpts carry no citation command at all, so no bibliography entry of the article is transcribed and no reference is redistributed. Reference adaptation: none was needed and none was made; every reference inside the delivered unit resolves inside it, so no unresolved reference or citation is produced. Printed slips kept literal: no printed word, symbol, hypothesis or number of the counted statement or of its proof was altered. Layout and class: the article's own class is \texttt{amsart} with packages including graphicx, color, amsmath, amssymb, amsthm, amscd, amsfonts, float, enumerate, lineno, blindtext, subcaption, hyperref, fancyhdr; the wrapper is the lane's pinned \texttt{amsart} editorial wrapper at 10pt on A4, which restates the theorem-like environments the retained text needs and adds the article's own macro definitions that the retained text needs (none), with extra wrapper packages (bm, bbm, dsfont, xspace, cancel, mathtools). The delivered numbering is the unit's own. Assets: no retained span contains a figure environment, a graphics-inclusion command, a PDF-page inclusion, a table or a TikZ picture; no image file is copied into this package, the package declares no asset dependency and no image file is redistributed with it. Licence: version arXiv:2509.25464v1 of this article is distributed under the Creative Commons Attribution 4.0 International licence, as recorded in the arXiv licence metadata for this exact version and shown on the abstract page https://arxiv.org/abs/2509.25464v1. A full rights notice is delivered at licenses/arxiv-2509.25464-CC-BY-4.0.txt. Characters: every retained excerpt is pure ASCII, so the delivered engine, XeLaTeX with its default Latin Modern text font, reports no missing character.
\begin{thm}
    Given 2 vertices and k edges, there are $\frac{n(n+1)(3k-4n+1)}{3}+(n+1)\lceil\frac{k+1}{2}\rceil$ Leavitt path algebras up to isomorphism where $n=\lceil\frac{k}{2}\rceil$.
\end{thm}

\begin{proof}
     As Leavitt path algebras are uniquely associated to directed graphs, this result is equivalent to the following, which we will prove: given 2 vertices and k edges, there are $\frac{n(n+1)(3k-4n+1)}{3}+(n+1)\lceil\frac{k+1}{2}\rceil$ graphs up to isomorphism, where $n=\lceil\frac{k}{2}\rceil$. To prove this, we begin by fixing the number of loops on each vertex (call them $u$ and $v$) and counting the possible arrangements of edges between them. Furthermore, by symmetry, we fix the number of loops on $u$ to be greater than or equal to the number of loops on $v$. With this in mind, we can start counting the possibilities given some $k$. 
     
    If there are 0 loops on $u$, then $\lceil\frac{k+1}{2}\rceil$ graphs are added. Namely, each of the $k$ vertices can go from $u$ to $v$ or $v$ to $u$, yielding $k+1$ possibilities, but symmetry requires those possibilities to be divided by 2. However, if $k$ is even, then $\frac{k+1}{2}\notin\mathbb{Z}$. Considering more concretely what takes place with an even number of edges, we see that there are $\frac{k}{2} + 1$ different graphs. These graphs are constructed by starting with the graph with all edges facing one direction, then flipping one edge in the opposite direction at a time until half of the edges are facing the opposite direction. Noticing that for odd $k$ we have $\frac{k+1}{2}$ different graphs and for even $k$ we have $\frac{k}{2} + 1 = \frac{k+2}{2}$ different graphs, we see that we have $\lceil\frac{k+1}{2}\rceil$, as claimed. 
    
    If there is 1 loop on $u$, then we have $k$ possible graphs where there are 0 loops on $v$ made from flipping the $k-1$ edges between $u$ and $v$ one at a time, and by a similar argument as the case with 0 loops on $u$, if we have one loop on $u$ and $v$, we get $\lceil\frac{k-1}{2}\rceil$. Hence, we have $k+\lceil\frac{k-1}{2}\rceil$ different graphs.
    
    More generally, we see that given $m$ loops on $u$, we have $(k-m+1) + (k-m) + (k-m-1) +\dots + (k-m-(m-1)) + \lceil\frac{k-(2m-1)}{2}\rceil$ different graphs, where each term adds all graphs with one more loop added to $v$ and the final ceiling term adds the graphs where $v$ also has $m$ loops.

    The sum of all of these possible graphs can now be expressed as the sum of three different sums:
    \begin{itemize}
        \item [($i$)] $k + 2(k-2) + 3(k-4) + \dots $
        \item [($ii$)] $(k-1) + 2(k-3) + 3(k-5) + \dots $
        \item [($iii$)] $\lceil\frac{k+1}{2}\rceil + \lceil\frac{k-1}{2}\rceil + \lceil\frac{k-3}{2}\rceil + \dots $
    \end{itemize}
    However, notice that for fixed $k$, these sums need to be terminated such that negative terms are not counted. So, we must define the sums in such a way that $(k-c)$ terms are not added for any $c>k$. \\
    We begin by combining sums ($i$) and ($ii$). Rearranging terms, we get that together they form the sum $2(1+2+3+\dots )k - ((0+1) + 2(2+3) + 3(4+5) + \dots )$, which has general term $2nk - n(4n-3)$. Notice that this general term breaks into the two terms $n(k - (2n-2))$ and $n(k-(2n-1))$, terms from ($i$) and ($ii$), respectively. We want to terminate this series at an $n$ such that all positive $(k-(2n-2))$ and $(k-(2n-1))$ are included while no negatives are. In other words, we want to find the appropriate $n$ for the series $\sum_{i=1}^n (2ik - i(4i-3))$. \\
    
    \textbf{Claim:} $n = \lceil\frac{k}{2}\rceil$. \\

    To demonstrate this claim, we consider 2 cases: (1) $k$ is even, (2) $k$ is odd. \\
    First, suppose $k$ is even. Then, $n = \frac{k}{2}$, which implies $k = 2n$. Now, we consider the last terms in the sum, namely $n(k-(2n-2))$ and $n(k-(2n-1))$. We see that $n(k-(2n-2)) = 2n\geq 0$ and $n(k-(2n-1)) = n\geq 0$. If we added another term to the series, we would have the terms $(n+1)(k-(2(n+1)-2)) = 0$ and $(n+1)(k-(2(n+1)-1)) = -(n+1) < 0$, which we clearly do not want to include. Hence, for even $k$, we see that $n = \lceil\frac{k}{2}\rceil$ is the appropriate stopping point. \\
    Now, suppose $k$ is odd. Then, $n = \frac{k+1}{2}$, which implies $k = 2n-1$. Once again, we consider the last terms, namely $n(k-(2n-2))$ and $n(k-(2n-1))$. We see that $n(k-(2n-2)) = n\geq 0$ and $n(k-(2n-1)) = 0\geq 0$. If we added another term to the series, we would have the terms $(n+1)(k-(2(n+1)-2)) = -(n+1) < 0$ and $(n+1)(k-(2(n+1)-1)) = -2(n+1) < 0$, which we do not want to include. Hence, for odd $k$, we see that $n = \lceil\frac{k}{2}\rceil$ is the appropriate stopping point as well, proving our claim. \\

    To summarize, combining sums ($i$) and ($ii$) and accounting for when they should stop, we have $\sum_{i=1}^n (2ik-i(4i-3))$, where $n = \lceil\frac{k}{2}\rceil$. Notice that starting the sum at $i=0$ has no effect on the sum, so we will rewrite our sum as $\sum_{i=0}^n 2ik-i(4i-3)$, where $n = \lceil\frac{k}{2}\rceil$. \\
    Now, we consider ($iii$), namely the sum $\lceil\frac{k+1}{2}\rceil + \lceil\frac{k-1}{2}\rceil + \lceil\frac{k-3}{2}\rceil + \dots $. This sum has general term $\lceil\frac{k-(2n-1)}{2}\rceil$, so we have the sum $\sum_{i=0}^n\lceil\frac{k-(2i-1)}{2}\rceil$. Once again, we need to find the appropriate $n$ such that no negative terms are included. \\

    \textbf{Claim:} $n = \lceil\frac{k}{2}\rceil$. \\

    To see this, we study the cases when $k$ is even and $k$ is odd once again. Supposing $k$ is even, we have $n = \frac{k}{2}$, which implies $k = 2n$. With this, we see that $\lceil\frac{k-(2n-1)}{2}\rceil = \lceil\frac{1}{2}\rceil = 1\geq 0$, while the next term $\lceil\frac{k-(2(n+1)-1)}{2}\rceil = \lceil\frac{-1}{2}\rceil = 0$. Seeing that the term after our stopping point is 0 and that subsequent terms would be negative, we see that $n = \lceil\frac{k}{2}\rceil$ is an appropriate stopping point for even $k$. Supposing $k$ is odd, we have that $n = \frac{k+1}{2}$, which implies $k = 2n-1$. Considering the last term in the sum, With this, we see that $\lceil\frac{k-(2n-1)}{2}\rceil = \lceil\frac{0}{2}\rceil = 0$, where any subsequent terms would be less than or equal to 0. With this, we see that $n = \lceil\frac{k}{2}\rceil$ is an appropriate stopping point for odd $k$, as desired. \\

    Hence, we now have that ($iii$) with the appropriate stopping point is $\sum_{i=0}^n\lceil\frac{k-(2i-1)}{2}\rceil$ where $n = \lceil\frac{k}{2}\rceil$. \\
    Noticing that both this sum and the combined sum of ($i$) and ($ii$) can be combined, we get the sum $\sum_{i=0}^n (2ik - i(4i-3) + \lceil\frac{k-(2i-1)}{2}\rceil)$ where $n = \lceil\frac{k}{2}\rceil$. We now manipulate this sum, which yields:

    \begin{align*}
        \sum_{i=0}^n \bigg(2ik - i(4i-3) + \bigg\lceil\frac{k-(2i-1)}{2}\bigg\rceil\bigg) &= \sum_{i=0}^n \bigg((2k+3)i - 4i^2 + \bigg\lceil\frac{k+1-2i}{2}\bigg\rceil\bigg)\\
        &= \frac{n(n+1)(2k+3)}{2} + \frac{2n(n+1)(2n+1)}{3}\\
        &\quad + \sum_{i=0}^n\bigg(\bigg\lceil\frac{k+1}{2}\bigg\rceil - i\bigg)\\
        &= \frac{n(n+1)(6k-8n+5)}{6} - \frac{n(n+1)}{2}\\
        &\quad + (n+1)\bigg\lceil\frac{k+1}{2}\bigg\rceil\\
        &= \frac{n(n+1)(3k-4n+1)}{3}+(n+1)\bigg\lceil\frac{k+1}{2}\bigg\rceil.
    \end{align*} 
\end{proof}

\end{document}
